% ============================================================
%  A2 - Testdokument für das eigene Paket mcstyle.sty
%  Kompilieren: zweimal pdflatex paket-test.tex
%  (zweiter Durchlauf wegen \gesamtpunkte)
% ============================================================
\documentclass[a4paper,11pt]{article}
\usepackage[T1]{fontenc}
\usepackage[ngerman]{babel}

% Varianten zum Ausprobieren:
%   \usepackage{mcstyle}            normale Ausgabe
%   \usepackage[entwurf]{mcstyle}   mit Entwurfsvermerk und Anmerkungen
%   \usepackage[sw]{mcstyle}        Schwarz-Weiß
\usepackage[entwurf]{mcstyle}

\begin{document}

\section*{Übungsblatt}
Erstellt an der \fh{}. Insgesamt sind \textbf{\gesamtpunkte{} Punkte} erreichbar.

\begin{hinweis}
  Alle Befehle auf dieser Seite stammen aus \texttt{mcstyle.sty}.
  Die Präambel bleibt dadurch kurz.
\end{hinweis}

\aufgabe[3]{Dokumentklassen}
Nennen Sie drei Standardklassen von \LaTeX. \anmerkung{Lösung ergänzen}
\antwortzeilen{3}

\aufgabe[1]{Kompilieren}
Welches Programm erzeugt aus einer \texttt{.tex}-Datei ein PDF?
\antwortzeilen{1}

\begin{hinweis}[Tipp]
  Der Kasten akzeptiert einen eigenen Titel als optionales Argument.
\end{hinweis}

\end{document}
